%====================================================================
% Chapter 2 — Background and Related Work
% Self-contained section file: no preamble, no \documentclass.
% Requires (provided by main.tex preamble): booktabs, tikz, amsmath,
% biblatex/natbib. TikZ libraries used: positioning, arrows.meta, fit.
%====================================================================
\chapter{Background and Related Work}
\label{ch:background}

This chapter follows the two-part recipe of a conventional literature review: a
concise technical primer that fixes the vocabulary used throughout the rest of
this document (Section~\ref{sec:bg-primer}), and a per-paper survey of the
feline and animal facial-pain literature that situates our work
(Section~\ref{sec:bg-related}). It closes with a synthesis
(Section~\ref{sec:bg-synthesis}) that contrasts prior work against the two
\emph{portable methods} this project headlines, and makes the
\enquote{not-a-replicate} argument explicit by enumerating the replication
traps a single-corpus cat-pain study must avoid.

A framing rule governs the whole chapter and is stated once here so it need not
be repeated: the modelling \emph{engine} of this project---a frozen DINOv2
backbone feeding rank-consistent ordinal heads---is treated as conceded
\emph{plumbing} and is never claimed as a novel contribution. The primer
therefore describes these components only to the depth needed to read the
method, not as from-scratch tutorials, and the synthesis is careful to locate
novelty exclusively in the two measurement protocols, not in the architecture.

\section{Technical background}
\label{sec:bg-primer}

\subsection{The Feline Grimace Scale and its action-unit rubric}
\label{subsec:bg-fgs}

The Feline Grimace Scale (FGS) is an observational pain instrument that scores a
cat's facial expression along five \emph{action units} (AUs): ear position,
orbital tightening, muzzle tension, whisker change, and head position. Each AU
is scored on a three-point ordinal scale ($0$~=~action unit absent,
$1$~=~moderately present or uncertain, $2$~=~obviously present), and the five
scores are summed and rescaled to a $0$--$10$ index; a normalised ratio
$\geq 0.39$ is the published decision boundary indicating that analgesia should
be considered \cite{evangelista2019}. The rubric wording adopted by the VLM
weak-labeller in this project is taken verbatim from this scale, so that the
labelling target is the same construct the validation study certified rather
than a paraphrase. Figure~\ref{fig:fgs-schematic} shows the AU-to-score
schematic that the rest of this document refers to.

Two properties of the scale matter for the methods that follow. First, the AU
scores are \emph{ordinal, not nominal}: a confusion of $0$ with $2$ is a larger
error than a confusion of $0$ with $1$, which is why every agreement and
calibration metric in this work is ordinal-aware rather than treating the three
levels as unordered classes. Second, the summed index is a
\emph{convolution of five low-cardinality variables}; the resulting $0$--$10$
sum has only eleven atoms, a fact that later forbids drawing a binned
reliability diagram over the sum (Section~\ref{subsec:bg-calibration}).

\begin{figure}[t]
  \centering
  \begin{tikzpicture}[
      font=\small,
      au/.style={draw, rounded corners, minimum width=2.4cm,
                 minimum height=0.8cm, align=center, fill=gray!8},
      lvl/.style={draw, circle, minimum size=0.55cm, inner sep=0pt,
                  font=\scriptsize},
      node distance=0.35cm and 1.4cm,
      >={Stealth[length=2mm]}
    ]
    % five AUs
    \node[au] (ear)  {Ear\\position};
    \node[au, below=of ear]    (orb) {Orbital\\tightening};
    \node[au, below=of orb]    (muz) {Muzzle\\tension};
    \node[au, below=of muz]    (whk) {Whisker\\change};
    \node[au, below=of whk]    (hed) {Head\\position};

    % per-AU ordinal levels
    \foreach \n/\y in {ear/0, orb/-1.15, muz/-2.3, whk/-3.45, hed/-4.6} {
      \node[lvl] (\n0) at (3.2,\y) {0};
      \node[lvl, right=0.25cm of \n0] (\n1) {1};
      \node[lvl, right=0.25cm of \n1] (\n2) {2};
      \draw[->] (\n.east) -- (\n0.west);
    }

    % aggregation node
    \node[au, fill=blue!8, minimum height=2.2cm, right=2.6cm of muz]
      (sum) {Sum and\\rescale\\$\sum_{i=1}^{5}\!\mathrm{AU}_i \to [0,10]$};
    \foreach \n in {ear,orb,muz,whk,hed} {
      \draw[->, gray] (\n2.east) -- (sum.west);
    }

    % decision
    \node[au, fill=red!8, right=1.4cm of sum] (dec)
      {Index $\geq 0.39$\\$\Rightarrow$ consider\\analgesia};
    \draw[->] (sum.east) -- (dec.west);
  \end{tikzpicture}
  \caption[FGS action-unit to score schematic]{The Feline Grimace Scale maps
    five facial action units, each scored ordinally on $\{0,1,2\}$, to a summed
    $0$--$10$ index, with a published decision threshold at the normalised ratio
    $0.39$ \cite{evangelista2019}. In this project the per-AU rubric is the
    weak-labelling target; summation and the $0.39$ flag are computed in code,
    never elicited from the rater.}
  \label{fig:fgs-schematic}
\end{figure}

\subsection{Ordinal regression and rank consistency}
\label{subsec:bg-ordinal}

Because each AU is an ordered three-level variable, predicting it with a plain
softmax classifier discards the ordering and admits rank-inconsistent
probability estimates (e.g.\ assigning high mass to $\Pr(\mathrm{AU}\geq 2)$
while assigning low mass to $\Pr(\mathrm{AU}\geq 1)$). The CORAL framework
\cite{cao2020coral} reformulates a $K$-level ordinal problem as $K-1$ binary
\enquote{is the label greater than rank $k$?} sub-tasks that share a single
feature projection plus per-rank biases, which guarantees the predicted
cumulative probabilities are monotone and hence rank-consistent. CORN
\cite{shi2023corn} relaxes the weight-sharing constraint of CORAL by conditioning
each rank's probability on the previous one through the chain rule, recovering
rank consistency without sacrificing the flexibility of independent rank
classifiers; it is the head we adopt for each AU. The relevant output is the set
of soft cumulative probabilities $\Pr(\mathrm{rank} > k)$, from which a per-AU
probability mass function over $\{0,1,2\}$ is read off. A captioned placeholder
for this head is shown in Figure~\ref{fig:corn-placeholder}; the convolution of
the five per-AU PMFs into a distribution over the $0$--$10$ sum is the
\emph{distributional decode} used for the calibration metrics below.

\begin{figure}[t]
  \centering
  \fbox{\parbox{0.86\linewidth}{\centering\vspace{0.4cm}
    \textbf{[Schematic placeholder --- to be produced]}\\[0.3cm]
    \small CORN rank-consistent ordinal head: a frozen feature vector feeds a
    shallow projection producing $K-1=2$ conditional rank logits per action
    unit; $\Pr(\mathrm{rank}>0)$ and $\Pr(\mathrm{rank}>1)$ are chained into a
    monotone cumulative profile and decoded to a PMF over $\{0,1,2\}$. Five such
    heads are convolved into a PMF over the $0$--$10$ sum.\vspace{0.4cm}}}
  \caption[CORN ordinal head schematic (placeholder)]{Captioned placeholder for
    the per-AU CORN head \cite{shi2023corn,cao2020coral}. No results are
    depicted; this is a structural diagram to be rendered in the methods
    chapter.}
  \label{fig:corn-placeholder}
\end{figure}

\subsection{Self-supervised frozen vision features}
\label{subsec:bg-dinov2}

DINOv2 \cite{oquab2023dinov2} is a self-supervised vision transformer trained
without labels on a large curated image corpus, producing general-purpose
features that transfer to downstream tasks without fine-tuning. We use a frozen
DINOv2 ViT-S/14 encoder purely as a fixed feature extractor: the backbone is
never updated, its pooled features are cached to disk once, and only the shallow
CORN heads are trained on the cache. This is the appropriate low-data choice on
the available hardware---fully fine-tuning a transformer on a few hundred faces
would overfit catastrophically---and it is deliberately unremarkable. The
backbone is part of the conceded engine; we make no claim that the choice of
self-supervised features is itself a contribution.

\subsection{Vision--language models as structured raters}
\label{subsec:bg-vlm}

A vision--language model (VLM) can be prompted to read an image and emit a
structured answer. We use this capability narrowly: the VLM is constrained by a
forced tool-use schema to return five ordinal AU scores, each a closed
enumeration over $\{0,1,2\}$, with a rationale produced before each score. The
VLM is therefore a \emph{weak labeller} that proposes per-AU FGS scores at
scale and cheaply, which a veterinarian then accepts or corrects. The central
empirical question this raises---and the first headline of this work---is
whether such a frozen VLM can weak-label feline FGS action units at
human-rater agreement, measured against a veterinary anchor. That measurement,
not the act of prompting a VLM, is the contribution; the rater is selected by
its measured agreement rather than asserted to be adequate.

\subsection{Calibration and the expected calibration error}
\label{subsec:bg-calibration}

A model is \emph{calibrated} when its stated probabilities match empirical
frequencies. The expected calibration error (ECE) summarises miscalibration by
binning predictions and comparing confidence to accuracy within each bin;
temperature scaling \cite{guo2017calibration} is a one-parameter post-hoc fix
that rescales logits to reduce ECE without changing the ranking of predictions.
For the ordinal AUs we report a \emph{per-AU classwise} ECE so that
miscalibration is attributable to a specific action unit, and we calibrate each
CORN head with a single temperature scalar fit on training folds. A deliberate
restriction follows from Section~\ref{subsec:bg-fgs}: because the $0$--$10$ sum
has only eleven atoms, a binned reliability diagram on the sum is degenerate at
this sample size, so the distributional path reports a ranked probability score
on the sum plus per-AU classwise ECE, and \emph{no} binned reliability diagram
on the sum is ever drawn. Calibration is supporting hygiene in this work, not a
headline.

\subsection{Decision curve analysis and net benefit}
\label{subsec:bg-dca}

Generic accuracy and calibration metrics ignore that the two error types of a
welfare instrument are not equally costly: failing to flag a painful cat is far
worse than over-flagging a comfortable one. Decision curve analysis
\cite{vickers2006dca} addresses this by plotting \emph{net benefit} against a
threshold probability that encodes the clinician's relative weighting of the two
errors, allowing a model to be judged across a \emph{range} of harm ratios
rather than at a single symmetric-cost operating point. Because we have no
veterinarian-elicited undertreat:overtreat ratio, we sweep the ratio as a range
along the threshold-probability axis. This welfare-asymmetric decision curve is
the headline artifact of the trustworthiness wrapper---a clinical-decision
contribution that distinguishes the wrapper from generic calibration---and the
operating point itself is selected at a fixed pain-recall $\geq 0.90$ anchored
to the FGS validation study, never at a symmetric Youden-$J$ or $F_1$ knee.

\subsection{Conformal risk control: Learn-Then-Test}
\label{subsec:bg-ltt}

Learn-Then-Test (LTT) \cite{angelopoulos2021ltt} is a distribution-free,
finite-sample procedure for choosing model thresholds so that a user-specified
risk is controlled with high probability. We use it for the defer-to-vet
abstention rule: two thresholds bracket the $0.39$ point so that confident
negatives and confident positives are returned while the uncertain middle is
deferred, and LTT certifies a one-sided lower bound on the negative predictive
value (NPV, i.e.\ control of missed pain) at each abstention rate. The output is
reported strictly as a \emph{lower bound}; the word \enquote{guaranteed} is
banned, and if the available label budget cannot certify a useful band the
abstention curve is reported as exploratory rather than as a controlled claim.

\subsection{Inter-rater agreement: weighted kappa and Krippendorff's alpha}
\label{subsec:bg-agreement}

Agreement between two raters on an ordinal scale is measured with the
quadratically weighted kappa \cite{cohen1968kappa}, which penalises
disagreements by the squared distance between levels and is therefore the
correct statistic for the three-level AUs (a nominal kappa over
$\{0,1,2\}$ without weights is a bug to avoid). The headline measurement of this
project is exactly such a per-AU quadratic weighted kappa between the VLM and
the veterinary anchor, gated on its confidence-interval \emph{lower bound}
rather than the point estimate, because at the small anchor sizes available a
point estimate of $0.6$ is statistically indistinguishable from a true value of
$0.47$. Separately, Krippendorff's alpha \cite{krippendorff2004} in its ordinal
form measures the self-consistency of the VLM across repeated runs at non-zero
temperature; this is a vet-free reliability pre-screen and an abstention signal,
and it complements but never substitutes for the vs-vet kappa.

\subsection{Confident learning for label triage}
\label{subsec:bg-confident}

Confident learning \cite{cleanlab2021} estimates which labels in a noisy dataset
are likely wrong by analysing the joint distribution of given and predicted
labels. We use it not to clean labels automatically but to \emph{triage} the
scarce veterinary review budget: images that confident learning flags as likely
mislabelled, together with VLM disagreements and near-threshold cases, are sent
to the veterinarian first, maximising label corrections per review minute.

\subsection{Background robustness and saliency faithfulness}
\label{subsec:bg-robustness}

A face-pain classifier can appear accurate while reading the wrong signal---
lighting, background, or capture conditions correlated with the pain label
rather than the face itself. The backgrounds-challenge line of work
\cite{xiao2021background} demonstrates that models routinely exploit such
spurious foreground/background cues and provides the conceptual basis for our
\emph{capture-condition confound audit}: a trivial classifier built on
brightness, blur, aspect ratio, and generic embedding features is trained to
predict pain, and if it beats chance the label is entangled with acquisition
context. Complementarily, the energy-based pointing game from the Score-CAM line
\cite{wang2020scorecam} quantifies saliency \emph{faithfulness}---how much of a
saliency map's energy falls inside the relevant region---and underlies the
per-AU energy-based pointing game (EBPG) we report alongside the
foreground/background gap (FGS-BG-Gap). Crucially this audit is
\emph{one-directional}: it is well powered to \emph{detect} confounding but
underpowered to rule it out, so its only honest conclusion is \enquote{no
confound detected at this power,} never \enquote{no confound.} The
\emph{protocol}---FGS-BG-Gap plus per-AU EBPG, runnable on any future
facial-pain corpus---is the second headline contribution, distinct from any
verdict about a particular dataset.

\section{Related feline and animal pain work}
\label{sec:bg-related}

Each entry below follows the same recipe: aim, method, the key reliability
statistic the authors report, and the limitation that bears on our design. The
reported numbers belong to the cited studies and are quoted here as background;
this project has not yet run experiments and reports no results of its own.

\subsection{FGS validation (Evangelista et al.)}
The aim was to construct and validate the Feline Grimace Scale as an
observational acute-pain instrument \cite{evangelista2019}. The method paired
human raters scoring the five AUs against established pain references, with
psychometric validation of reliability, criterion validity, and a decision
threshold. The headline reliability statistics are an area under the curve of
approximately $0.94$ with sensitivity around $90.7\%$ and specificity around
$86.6\%$ at the normalised-ratio threshold of $0.39$. The limitation for us is
that this is a \emph{human-rater} instrument: it certifies that trained people
scoring the rubric recover pain, not that any automated pipeline does, and it
fixes the construct and the operating anchor (pain-recall $\geq 0.90$) we
inherit rather than a model to compare against.

\subsection{Automated landmark FGS (Steagall et al.)}
The aim was to automate FGS scoring from images \cite{steagall2023}. The method
detected facial landmarks, derived geometric features from them, and classified
with a gradient-boosted model (XGBoost), with a minimum-aggregation rule across
whisker and head regions. The reported headline accuracy is about $95.5\%$. Two
limitations shape our positioning: the pipeline is a \emph{landmark-geometry}
approach that we deliberately do not replicate (it is one of the replication
traps in Section~\ref{sec:bg-synthesis}), and a single bare accuracy figure on
its own corpus is precisely the anti-benchmark we refuse to compete against on
headline terms.

\subsection{Explainable binary cat-pain detection (Feighelstein et al.)}
The aim was binary cat-pain detection with explanation \cite{feighelstein2023}.
The method compared a landmark-based classifier (reported around $77\%$) against
a deep-learning classifier (reported around $65\%$), with an explainability
analysis indicating the mouth region contributed more than the eyes, which
contributed more than the ears, and a two-source pain ground-truth rule
(a composite pain-scale threshold together with a charted clinical reason). The
limitation is that this is a \emph{binary} detector on a related corpus; without
the firewall and per-cat grouping we adopt, the natural next paper is a
backbone-swapped re-run of exactly this, which is the study we are explicitly
trying not to write.

\subsection{CatFLW landmarks (Martvel et al.)}
The aim was a cat facial-landmark dataset and detector \cite{martvel2024catflw}.
The method provided CatFACS-aligned facial landmarks with reported normalised
mean error in roughly the $9$--$26\%$ range at a modest landmark count. The
limitation, and the reason it is preprocessing rather than an anchor in our
pipeline, is that it carries landmarks and bounding boxes but \emph{no pain or
AU labels}; treating it as a graded FGS anchor would be a category error.
We use it only for an eye-alignment / NME audit to confirm crops are
geometrically clean enough to feed the ordinal heads.

\subsection{COSMIN scoping review of human-rater instruments (Lee \& Steagall)}
The aim was a COSMIN-style scoping review of the measurement properties of
feline pain instruments \cite{leesteagall2026}. The method screened and
appraised \emph{human-rater} pain-assessment instruments against measurement-
property criteria. The limitation that governs how we cite it is one of scope:
the review covers human-rater instruments only, and automated/AI scoring is
outside its topical remit. We therefore do not claim the review named
calibration or confidence intervals as gaps, nor that it excluded automated
scorers as a finding---it simply did not address them.

\subsection{Low-prevalence measurement-error propagation (Chavoshi et al.)}
The aim was to characterise how measurement error propagates at low disease
prevalence \cite{chavoshi2025}. The method analysed how imperfect labels and
classifiers degrade effective performance as prevalence falls, with a worked
illustration in which a nominally strong classifier collapses to roughly $53\%$
effective sensitivity around $10\%$ prevalence. The implication for us is direct
and sobering: our corpus sits near $12$--$13\%$ pain prevalence, so headline
accuracy is the wrong target, confidence intervals and a fixed-high-sensitivity
operating point are mandatory, and any single-number boast would be misleading
in exactly the regime this paper formalises.

\section{Synthesis and positioning}
\label{sec:bg-synthesis}

\subsection{What is new here, and what is conceded}
The prior work above clusters into two families: detectors that report a single
accuracy figure on their own corpus
\cite{steagall2023,feighelstein2023}, and the instrument/measurement context
that frames them \cite{evangelista2019,leesteagall2026,chavoshi2025} plus the
landmark scaffolding \cite{martvel2024catflw}. Our contribution is deliberately
\emph{not} a better detector. It is two transportable measurement protocols that
no prior feline-pain study produced: (i) a measurement of whether a frozen VLM
can weak-label FGS action units at human-rater agreement, expressed as five
per-AU quadratic weighted kappas against a veterinary anchor and gated on the
confidence-interval lower bound; and (ii) a reusable capture-condition
confound-attribution protocol (FGS-BG-Gap plus per-AU EBPG), one-directional by
construction. The DINOv2-plus-CORN engine carrying these methods is conceded as
plumbing and is never claimed as novel; the binary-plus-wrapper spine ships with
a welfare-asymmetric decision curve and a one-sided NPV lower-bound abstention
curve, while the $0$--$10$ ordinal layer ships inspected-but-not-validated.
Table~\ref{tab:positioning} contrasts the prior work with this positioning, with
accuracy treated explicitly as an anti-benchmark and a column marking whether a
contribution is a \emph{portable method}.

\begin{table}[t]
  \centering
  \caption[Prior work versus this work positioning matrix]{Positioning matrix.
    Accuracy figures are quoted from the cited studies as background and are
    treated as an \emph{anti-benchmark}: this work does not compete on a single
    corpus accuracy. The final column marks whether the contribution is a
    dataset-agnostic \emph{portable method}. Quoted statistics are the original
    authors'; this project reports no results.}
  \label{tab:positioning}
  \small
  \begin{tabular}{@{}p{3.0cm}p{3.1cm}p{3.3cm}cc@{}}
    \toprule
    \textbf{Work} & \textbf{Method} & \textbf{Reported reliability stat
      (theirs)} & \textbf{Task} & \textbf{Portable method?} \\
    \midrule
    Evangelista et al.\ \cite{evangelista2019}
      & Human-rater FGS, validated
      & AUC $\approx 0.94$; sens.\ $90.7\%$; spec.\ $86.6\%$ at $0.39$
      & Graded (human) & No (instrument) \\
    Steagall et al.\ \cite{steagall2023}
      & Landmark geometry $\to$ XGBoost
      & Acc.\ $\approx 95.5\%$
      & Graded (auto) & No \\
    Feighelstein et al.\ \cite{feighelstein2023}
      & Landmark vs.\ deep, with XAI
      & Acc.\ $\approx 77\%$ (landmark) / $65\%$ (DL)
      & Binary & No \\
    Martvel et al.\ \cite{martvel2024catflw}
      & CatFACS landmark detector
      & NME $\approx 9$--$26\%$
      & Landmarks & No (no pain labels) \\
    Lee \& Steagall \cite{leesteagall2026}
      & COSMIN scoping review
      & Measurement-property appraisal (human raters only)
      & Review & No \\
    Chavoshi et al.\ \cite{chavoshi2025}
      & Error-propagation analysis
      & $\sim 53\%$ eff.\ sens.\ at $10\%$ prevalence
      & Theory & Yes (transfers) \\
    \midrule
    \textbf{This work (i)}
      & VLM-as-AU-rater, vs.-vet kappa
      & \emph{CI-lower-bound-gated per-AU QWK (to be measured)}
      & AU rating & \textbf{Yes} \\
    \textbf{This work (ii)}
      & Confound-attribution protocol
      & \emph{FGS-BG-Gap + per-AU EBPG (one-directional)}
      & Audit & \textbf{Yes} \\
    \textbf{This work (frame)}
      & Welfare-asymmetric DCA + LB-abstention
      & \emph{net benefit (range); one-sided 95\% NPV LB}
      & Binary wrapper & \textbf{Yes} \\
    \bottomrule
  \end{tabular}
\end{table}

\subsection{The not-a-replicate case and eight replication traps}
\label{subsec:bg-traps}

Because this project is instantiated on a single in-domain corpus, the most
likely way it could collapse into \enquote{another cat-pain detector} is by
falling into one of a small set of replication traps. We name them here so the
methods chapter can be read as a list of deliberate avoidances:

\begin{enumerate}
  \item \textbf{Landmark-geometry replication.} Re-implementing the Steagall
    landmark-to-XGBoost pipeline \cite{steagall2023}; we use frozen features and
    ordinal heads precisely to not retread it.
  \item \textbf{Binary-detector replication.} Producing a backbone-swapped
    re-run of the Feighelstein binary detector \cite{feighelstein2023}; the
    portable methods, not the detector, are the headline.
  \item \textbf{Video/temporal replication.} Borrowing temporal landmark-track
    machinery; out of scope and not what the corpus supports.
  \item \textbf{External-cohort overclaim.} Asserting cross-cohort
    generalisation from a corpus dominated by a single known camera identity.
  \item \textbf{Single-rater ICC overclaim.} Reporting multi-rater agreement
    statistics when only one veterinary anchor exists.
  \item \textbf{Accuracy-as-benchmark.} Writing \enquote{we beat $77/79/95\%$};
    we always print the distinct-pain-cat denominator with Clopper--Pearson or
    bootstrap intervals, and never compete on a single corpus accuracy.
  \item \textbf{Cross-species headline.} Carrying a horse-grimace number into
    the abstract, results, or any transfer table; the horse set is decode
    scaffolding for a unit test only.
  \item \textbf{Validated-graded overclaim.} Reporting the $0$--$10$ graded
    layer as a validated instrument; it ships inspected-but-not-validated, and
    the word \enquote{graded} is struck from every validated-claim sentence.
\end{enumerate}

Avoiding these eight traps is what converts a single-corpus study into a
contribution that survives the obvious \enquote{swap the backbone and add
metrics} critique: the two surviving claims are facts about the world (a
measured VLM-vs-vet agreement) and a reusable audit protocol, both portable to
the next dataset, riding on an explicitly conceded engine.
